\scanpage{92}
rappelons-le, équivaut à la donnée de $(E,q,\omega)$, ou encore à celle de $\mathcal M$. Introduisons
\[
S[\underline i]=\operatorname{Spec}\underbrace{\underline O_S[\underline i]}_{=\underline O_S[T]/(T^2+1)}\tag*{(13)}
\]

On va définir une immersion ouverte dominante
\[
\Sigma\hookrightarrow X^{S[\underline i]}=\underline{\operatorname{Hom}}_S(S[\underline i],X)\tag*{(14)}
\]

Pour ceci, \deleted{explicitons les sections de $X^{S[\underline i]}$ sur $S$ comme \unsure{quotient|préliminaires}} \inserted{notons}
\[
\left\{\begin{aligned}&X\simeq\widetilde X/G_m\\&\widetilde X\subset\check{\mathbb V}(E)^*\qquad\widetilde X\ \text{d'équation}\quad q(x)=0\end{aligned}\right.\tag*{(15)}
\]

(NB $\widetilde X$ est un torseur sur $X$ de Groupe $G_{m,X}$, c'est le torseur associé au Module $\underline\Omega^1_{X/S}$ de degré relatif $2$\ldots)

\deleted{\unsure{On|Dire}} d'où on conclut que $\widetilde X^{\deleted{\unsure{[i]}}S[\underline i]}$ est un $G_m^{S[\underline i]}$-torseur sur $X^{S[\underline i]}$. Or les sections \inserted{$\unsure{\zeta|\xi}$} de
$\widetilde X^{S[\underline i]}\unsure{/|\big/}$ \deleted{sur} $S$ s'explicitent comme des sections de $\check{\mathbb V}(E)^*_{S[\underline i]}\inserted{{}\subset\mathbb V(E)_{S[\underline i]}}$ i.e. \deleted{des couples} comme
\[
\zeta=x+\underline i y\qquad x,y\in\Gamma(S,E),
\]
et il faut écrire
\[
\begin{aligned}\text{a)}\qquad &x+\underline i y\ \deleted{\text{ne s'annulent \unsure{simultanément} en aucune fibre}}\ \inserted{\text{comme section de }E_{S[\underline i]}\text{ ne}}\\&\text{s'annule en aucune fibre}\quad(\Rightarrow x\ \text{et }y\\&\text{ne s'annulent simultanément en aucune fibre}\end{aligned}\tag*{(16)}
\]
\unsure{bien sûr|est-ce sûr}\ldots)
% archive sheet number: 91
